\section{Eliminating the payment rule}
\label{sec:elimination}

The feasibility question of Section~\ref{sec:model} couples two unknowns: the
record $\enc$ and the payment decoders $\pay_i$. Products of the two make the
natural formulation bilinear, and a direct search over both does not scale. Our
first result removes $\pay_i$ entirely.

\begin{theorem}[Payment-decoder elimination]
\label{thm:elim}
Fix a record $\enc$ and an agent $i$. There exists a bounded decoder
$\pay_i \in [0,1]^K$ satisfying \eqref{eq:margin} if and only if
\[
  \min_{\lambda \in \simplex(\{0,1\}^{n-1})}
  \tfrac{1}{2}\bigl\|\,\E_{\lambda}\bigl[d_i\bigr]\,\bigr\|_1 \;\ge\; \gamma ,
\]
where $\E_\lambda[d_i] = \sum_{\sminus} \lambda(\sminus)\, d_i(\sminus)$.
\end{theorem}

\begin{proof}
Feasibility of \eqref{eq:margin} for agent $i$ says exactly
\[
  \max_{\pay \in [0,1]^K} \; \min_{\sminus} \; \langle \pay, d_i(\sminus)\rangle
  \;\ge\; \gamma .
\]
The inner minimum is over a finite set, so it equals the minimum over its convex
hull, i.e.\ over $\lambda \in \simplex(\{0,1\}^{n-1})$. The map
$(\pay,\lambda) \mapsto \langle \pay, \E_\lambda[d_i]\rangle$ is bilinear and both
$[0,1]^K$ and the simplex are convex and compact, so von Neumann's minimax
theorem applies:
\[
  \max_{\pay} \min_{\lambda} \langle \pay, \E_\lambda[d_i]\rangle
  \;=\; \min_{\lambda} \max_{\pay} \langle \pay, \E_\lambda[d_i]\rangle
  \;=\; \min_{\lambda} \sum_y \bigl(\E_\lambda[d_i]\bigr)_y^{+},
\]
the last step because maximising a linear form over the unit cube puts
$\pay_y = 1$ on positive coordinates and $\pay_y = 0$ elsewhere. Finally each
$d_i(\sminus)$ is a difference of two probability vectors, so
$\sum_y d_i(\sminus)_y = 0$, and the same holds for any mixture; for a
zero-sum vector $v$ we have $\sum_y v_y^+ = \tfrac12\|v\|_1$.
\end{proof}

\paragraph{What is and is not new here.}
The two ingredients are classical: von Neumann's minimax theorem, and the
variational characterisation of total variation as the largest gap a bounded
test function can produce. The \emph{qualitative} content---that holding an
agent to account requires its deviation to move the distribution of the observed
signal---is also standard, and is exactly the role played by the individual and
pairwise full-rank conditions of Fudenberg, Levine and
Maskin~\cite{flm1994} in repeated games with imperfect public monitoring. We
claim neither. What Theorem~\ref{thm:elim} adds is a \emph{quantitative,
one-shot} form of that condition---the shift must be at least $\gamma$ in total
variation, robustly over mixtures of teammate behaviour---and, decisively, the
elimination of the payment rule from the problem. Detectability conditions in the
repeated-games literature are stated for a \emph{given} monitoring structure;
here the monitoring structure is the unknown, and removing $\pay$ is what makes
it possible to ask which structures a budget can buy.

Three things follow, and the third is the one to remember.

\begin{corollary}[Edge condition]
\label{cor:edge}
Taking $\lambda$ a point mass gives a necessary condition: for every agent $i$
and every teammate profile $\sminus$,
\[
  \TV\bigl(\enc[(1,\sminus),\cdot],\; \enc[(0,\sminus),\cdot]\bigr)
  \;\ge\; \gamma .
\]
\end{corollary}

\noindent\textbf{(i) The problem becomes an optimisation over the record alone.}
Theorem~\ref{thm:elim} replaces a bilinear search over $(\enc, \pay)$ by a
condition on $\enc$. Computationally this is what makes
Section~\ref{sec:threshold} reachable: a bilinear formulation did not terminate
at $n = 7$, whereas the reduced question is a feasibility problem over $2^n$
binary variables and runs in milliseconds.

\noindent\textbf{(ii) Impossibility results become cheap.}
If a record maps some profile and all of its $n$ unilateral neighbours to the
same label, every edge total variation at that profile is $0$, so by
Corollary~\ref{cor:edge} no bounded payment achieves any positive margin---for
any decoder, not merely a binary one.

\noindent\textbf{(iii) What accountability \emph{is}.}
Corollary~\ref{cor:edge} is the general statement of the phenomenon this paper
is about:
\begin{quote}
\emph{A record supports accountability exactly to the extent that each agent,
acting alone, visibly moves it.}
\end{quote}
Reconstruction wants the opposite---it merges profiles that are cheap to confuse
in order to spend its $K$ labels well. The rest of the paper measures the
collision between those two demands.
